% 教学例 ex:fairness-rates 的计数；三面板保留比例分母，解释移至图注。
\begin{tikzpicture}[figure picture,foundation picture]
 \path[use as bounding box] (0,2) rectangle (168,47);
 \foreach \sx/\title in {0/人口统计均等,58/均等化赔率,116/正预测值均等}{
  \node at ({\sx+26},43) {\title};
 }
 \draw[draw=FigureMuted,line width=.5pt,densely dashed] (56,5)--(56,46);
 \draw[draw=FigureMuted,line width=.5pt,densely dashed] (114,5)--(114,46);
 % 人口统计均等：全体申请者为分母，条长对应24%与52%。
 \foreach \y/\group/\len/\ratio/\col/\fillcol in {34/u/9.6/24{\,/\,}100/FigureBlue/FigureBlueFill,18/v/20.8/52{\,/\,}100/FigureRed/FigureRedFill}{
  \node[anchor=east] at (5,\y) {$\group$};
  \fill[\fillcol] (8,{\y-2.5}) rectangle ({8+\len},{\y+2.5});
  \draw[draw=FigureStroke,line width=.7pt] (8,{\y-2.5}) rectangle (48,{\y+2.5});
  \draw[draw=\col,line width=.9pt] ({8+\len},{\y-2.5})--({8+\len},{\y+2.5});
  \node at (28,{\y-8}) {$\ratio$};
 }
 % 均等化赔率：给定资格后，两组均有TPR=.8、FPR=.1。
 \node[anchor=east] at (65,34) {$y=1$};
 \fill[FigurePanel] (66,31.5) rectangle (98,36.5);
 \draw[draw=FigureStroke,line width=.7pt] (66,31.5) rectangle (106,36.5);
 \draw[draw=FigureStroke,line width=.9pt] (98,31.5)--(98,36.5);
 \node at (86,26) {$16/20=48/60$};
 \node[anchor=east] at (65,18) {$y=0$};
 \fill[FigurePanel] (66,15.5) rectangle (70,20.5);
 \draw[draw=FigureStroke,line width=.7pt] (66,15.5) rectangle (106,20.5);
 \draw[draw=FigureStroke,line width=.9pt] (70,15.5)--(70,20.5);
 \node at (86,10) {$8/80=4/40$};
 % 正预测值均等：获选者为分母，条长对应16/24与48/52。
 \foreach \y/\group/\len/\ratio/\col/\fillcol in {34/u/26.6667/16{\,/\,}24/FigureBlue/FigureBlueFill,18/v/36.9231/48{\,/\,}52/FigureRed/FigureRedFill}{
  \node[anchor=east] at (121,\y) {$\group$};
  \fill[\fillcol] (124,{\y-2.5}) rectangle ({124+\len},{\y+2.5});
  \draw[draw=FigureStroke,line width=.7pt] (124,{\y-2.5}) rectangle (164,{\y+2.5});
  \draw[draw=\col,line width=.9pt] ({124+\len},{\y-2.5})--({124+\len},{\y+2.5});
  \node at (144,{\y-8}) {$\ratio$};
 }
\end{tikzpicture}%
